\section{평가 자료와 검토 절차}
\label{sec:data}

\begin{revision}

\subsection{평가 자료 개요}
\label{sec:data-overview}
핵심 평가는 설계자가 구성한 시나리오와 유한 범위의 전수 사건열을 사용한다. 자연 CoC는 사람에 따른 해석 차이와 기록 궤적의 대응 범위를 살피는 자료이며, 자연 오류 정확도를 계산하는 정답집으로 사용하지 않는다. 표~\ref{tab:evaluation-data}에 각 자료의 규모ㆍ작성 주체ㆍ역할과 결과 절을 정리했다.

\begin{table*}[t]
\centering
\color{revisionblue}
\bitablecaption{평가 자료의 구성과 용도. 각 행의 규모는 서로 다른 평가 단위다.}{Evaluation datasets, authorship and roles; counts across rows use different units.}
\label{tab:evaluation-data}
\footnotesize
\begin{tabularx}{\textwidth}{L{.18\textwidth}L{.13\textwidth}L{.16\textwidth}YL{.10\textwidth}}
\toprule
자료 & 규모 & 작성 주체 & 평가에서의 역할 & 결과 절 \\
\midrule
통제 시나리오 & 144건 & 설계자 & 검사기 비교ㆍ실패 분석ㆍ개발 후 같은 사례 재시험 & \ref{sec:first-results}--\ref{sec:agreement-results} \\
경계 시나리오 & 81건 & 설계자 & 장면 경계의 의무ㆍ증거 유지 효과 비교 & \ref{sec:boundary-results} \\
유한 전수 사건열 & 35,042,928건 & 생성기, 선별 없음 & 고정 구현과 별도 수식 판정기의 유한 범위 적합성 & \ref{sec:bounded-results} \\
자연 CoC 사람 검토 & 27장면ㆍ4인 & 자동 주석ㆍ사람 판정 & 주석 해석의 불일치와 궤적 대조의 한계 확인 & \ref{sec:human-results}ㆍ\ref{sec:trajectory-results} \\
자연 CoC 반응 분석 & 134사건 & 자동 주석ㆍ기록 속도 & 선택한 설정에서의 속도 반응 보조 분석 & \ref{sec:trajectory-results} \\
기존 회귀 자료 & 기준--변형 40쌍 & 설계자ㆍ변형 생성기 & 규칙 실행과 검사기 처리 오류의 회귀시험 & \ref{sec:regression-results} \\
\bottomrule
\end{tabularx}
\end{table*}

144건은 영어 원문과 구조화 입력을 함께 작성했고, 81건은 원 자연어 구절 없이 대상ㆍ조건ㆍ경계 연결을 제공한 자료다. 81건은 경계 확장 모델과 함께 설계하여 검증 실행 전에 고정했으며, 모델 고정 이후 새로 작성한 독립 시험으로 보지 않는다. 전수 사건열은 두 모델과 생성 규칙을 고정한 뒤 모든 선언 조합을 열거했다. 구체적인 조합과 비교 방법은 6장에서 설명한다.

\subsection{자연 자료의 출처와 선정}
\label{sec:provenance}
자연 자료는 nuScenes의 영상ㆍ자차 움직임과 사후 자동 생성한 CoC-Nusc 주석이다~\cite{caesar2020nuscenes,nvidia2026autolabeler}. 기록 궤적에서 계산한 행동 정보가 주석 생성에 쓰이므로 CoC와 궤적의 일치를 독립 정답으로 사용하지 않는다. 보관된 주석은 743장면ㆍ3,218사건이며, 이 중 자차 궤적 파일을 함께 보유한 94장면ㆍ403사건을 분석에 사용했다. CoC가 둘 이상인 90장면에서 만든 309인접쌍은 사람 검토의 후보 선정에 사용했으며 전체 쌍의 오류율을 산출하지 않았다.

CoC-Nusc의 배포 버전, 생성 실행별 모델 버전, 생성 지시문과 설정은 복원하지 못했다. 자료 설명서에서 모델 계열은 확인되지만 개별 문장의 정확한 생성 이력과 품질 점수의 계산식ㆍ통과 기준은 확인할 수 없다. 자료 선정에는 보관된 품질 통과 표시를 사용했다(표~\ref{tab:provenance}).

\begin{table}[t]
\centering
\color{revisionblue}
\bitablecaption{자료 출처의 확인 사실과 미복원 정보.}{Recovered provenance and information that remains unrecorded.}
\label{tab:provenance}
\footnotesize
\begin{tabularx}{\columnwidth}{L{.30\columnwidth}Y}
\toprule
항목 & 확인 범위 \\
\midrule
자연 입력 & 보관 파일ㆍ장면/사건 대응ㆍ파일 해시 \\
생성 모델 & Qwen/Qwen3.5-35B-A3B 계열 기재; 실행별 버전 미복원 \\
생성 지시문ㆍ설정 & 미복원 \\
품질 선정 & 저장된 통과 여부 사용; 점수 계산식ㆍ통과 기준 미복원 \\
\bottomrule
\end{tabularx}
\end{table}

보조 반응 분석에서는 403사건 중 품질 표시를 통과한 290건을 먼저 남겼다. 이 중 복합 행동이나 방향이 불명확한 156건을 제외한 134건을 사용했다. 134건은 속도 변화 탐지 대상 130건과 사건 시점에 이미 정지ㆍ양보 조건을 만족한 4건이다. 이 선정은 309인접쌍의 사람 검토 후보 선정과 별개다.

후보 선택기는 첫 문장에 정지ㆍ양보(\texttt{stop/yield}), 둘째 문장에 가속ㆍ진행(\texttt{accelerate/proceed}) 표현이 있는 쌍을 찾는다. 23쌍ㆍ15장면에서 장면별 가장 이른 창 하나를 선정했다. 비교군은 후보 장면을 제외하고 첫 문장에 정지ㆍ양보 표현이 있으나 둘째 문장에는 가속ㆍ진행 표현이 없는 창이다. 장면당 가장 이른 창을 남긴 뒤, 후보와의 사건 수 차이ㆍ정지/양보 종류 차이ㆍ원 사건 위치 차이 순으로 대응시켰다. 동률은 난수 시드 20260806과 창 ID의 해시로 결정하고 비교 창을 중복 사용하지 않아 12장면을 얻었다. 보존된 원 선정 프로그램을 다시 실행하면 기존 27개 검토 ID와 장면ㆍ집단 구분이 재현된다. 이 표본은 특정 행동 전환을 더 많이 포함하는 비무작위 표본이다.

309는 후보 선정에 사용한 인접 문장 쌍의 수이고, 27은 사람이 검토한 장면 묶음의 수다. 여러 사건이 한 장면에 속하므로 사건ㆍ쌍ㆍ장면을 독립 시행으로 합산하지 않는다.

\subsection{사람 검토 절차와 기록의 한계}
\label{sec:review-protocol}
보존된 검토 지침과 A/B 화면에 따르면, 1차 검토에는 CoC 원문과 첫 사건 기준 상대 시각을 제시하고 별도의 영상ㆍ궤적ㆍ객체/신호 상태 자료와 검사기 결과는 제공하지 않았다. 후보군ㆍ비교군 구분과 원 장면 ID는 화면에 표시하지 않고 집단 정보가 없는 검토 ID를 사용했다. 네 검토자는 시간 요구의 명시성, 요구 행동열과 문장 일관성을 판정했으며, 지침은 외부 객체 상태가 필요하면 미상으로 남기도록 했다. 원 판정과 최종 합의 기록은 보존되어 있으나 검토자 자격ㆍ교육 이력ㆍ독립 중재자 참여 및 당시 논의 절차는 확인할 기록이 부족하다. 합의 기록은 27장면 모두 텍스트 모순 없음이며 사유 메모는 0/27이다. 따라서 이 자료는 독립적인 자연 모순 정답이 아니라 주석 해석의 불일치와 모호성을 살피는 데 사용한다. 추가 사람 재평가는 수행하지 않았다.

재현을 위한 시나리오ㆍ생성 및 검사 코드ㆍUPPAAL 모델ㆍ질의ㆍ실행 기록을 보존했으며, 원 CoC 문장과 개별 검토 CSV는 접근 제한 자료로 관리한다.
\end{revision}
